% ===========================================================================
\section{Logical time in the runtime}\label{sec:time}
% ===========================================================================

\paragraph{Logical, not physical.}
All time in the theorems is \emph{model time}. A timestamp $t$ on an observation
means ``in the model, this event happens when the global logical clock $\theta$
equals $t$''. The runtime does not read the operating system clock in order to
decide semantic questions. The wall-clock time appears only as optional,
informational metadata in ledger records and is never consulted by the kernel.
The theorems therefore say nothing about how logical time relates to physical
time; establishing that relation (sensor timestamping, clock synchronisation,
bounded processing latency) is outside the scope of this report
(\cref{sec:limitations}).

\paragraph{Exact grid.}
Timestamps are parsed exactly from decimal strings (\code{twin::parse\_time}):
\code{"31.5"} at $R=1000$ is $31500$ ticks. A timestamp with more fractional
digits than the grid allows yields the error \code{TimeNotRepresentable}; it is
never rounded. The APIs reject JSON floating-point numbers for times.

\paragraph{Observation processing.}
An observation $(t,\sigma)$ is processed as a timed step: delay by $t-\theta$, then
the discrete step selected by $\sigma$ (\cref{def:monitor}). Hence:
\begin{itemize}
\item if $t<\theta$, the observation is rejected (\code{TimeRegression});
\item if the delay violates the current invariant, the observation is rejected
  (\code{IncompatibleObservation}, reporting the deadline \code{max\_delay}). For
  example, if the twin is in a location with invariant $x\le 5$ and $x=1.5$ at
  $\theta=28$, an event at $t=31.5$ requires $x=5$ and is admissible, whereas an
  event at $t=31.6$ is rejected;
\item a rejected observation leaves the semantic state unchanged and is recorded
  in the ledger as a rejection.
\end{itemize}

\paragraph{Simultaneous events.}
Events carrying the same timestamp are processed one after the other with zero
delay between them, in the order fixed by the runtime's single sequencer: the
order in which the input adapters enqueue them, with a deterministic
tie-breaking order between adapters. The ledger records that order, and replay
follows it. This is a total order on inputs; it is part of the input, not of
the model.

\paragraph{Pure passage of time.}
Between events the runtime may let time pass without an event
(\code{advance\_to}). This is the \textsc{(Delay)} rule. If the passage of time
alone would violate the current invariant (a missed deadline), the delay is
refused, the runtime reports a monitoring alarm, and the state is not changed.
